\section{Evaluation}
\label{sec:eval}

\subsection{Evaluation Setup}
We evaluate execution time, scaling and capacity, design contributions, numerical agreement, and parameter sensitivity under the following conditions.

\noindent\textbf{Hardware Specification.}
All experiments use NERSC Perlmutter GPU nodes, each with four NVIDIA A100 GPUs with 40~GB of device memory each, a 64-core CPU, and 256~GiB of host memory.
The nodes communicate over the Slingshot network.
GPU implementations use 90\% of available device memory, approximately 36~GB per GPU.
\SectorQ assigns one MPI rank to each node and one worker to each GPU, whereas distributed \cuPP assigns one rank to each GPU.
The largest configuration contains 128 nodes and 512 GPUs.
CPU implementations use the same nodes as the corresponding GPU configuration.
ORQA uses 64 ranks per node, MonoProp uses eight ranks with eight threads each, and threaded single-node implementations use 64 threads.
Qiskit uses one thread.
CPU and GPU comparisons use equal node counts, with the rank and thread configurations specified above.

\noindent\textbf{Evaluation Scope.}
The experiments measure Heisenberg-picture propagation that returns an observable expectation value.
Unless stated otherwise, each configuration is measured once with a 40-minute limit.
The size and exact-reference experiments use a 20-minute limit, and capacity ladders allow \SectorQ up to two hours per run.
A cutoff sweep stops after memory exhaustion or its time limit.
The plots distinguish completed measurements, observed failures, and unexecuted configurations.
Reported times cover propagation and any initialization performed lazily within the timed call, excluding queue and process-launch delays.
Horizontal-axis term counts denote the largest population at a \SectorQ window boundary under the stated truncation rule.
These counts differ from intermediate expansion within a window and the final output size.
Separate exact-reference and repeated-run experiments assess approximation error and timing variability.

\noindent\textbf{Circuits.}
The seven circuits in Table~\ref{tab:windows} vary the interaction graph, gate family, and observable.
The principal workload is a 127-qubit heavy-hex kicked Ising circuit with 20 steps and rotation angle $\pi/4$.
A second heavy-hex circuit uses transverse-field Ising model (TFIM) rotations at general angles.
Noisy kicked Ising adds depolarizing channels with probability 0.005 to the principal workload.
The remaining circuits are TFIM on an $11\times11$ lattice, XXZ on an $8\times8$ lattice, a 100-qubit QAOA circuit on a random regular graph, and a 100-qubit hardware-efficient ansatz.
Their respective depths are 10 steps, 6 steps, 4 layers, and 16 layers.
The observables include local magnetization and the edge sum of QAOA.
Additional TFIM grids and kicked Ising depths vary circuit dimensions.
Every comparison within a setting preserves the circuit, observable, and initial state.

\noindent\textbf{Baselines.}
The primary baseline is \cuPP~v0.5.0 with the original circuit and its default distributed mode.
The second baseline uses a Clifford-frame input and defers global deduplication until truncation, with additional merges where required to fit memory.
We report the fastest successful deferred interval among those tested.
The paulistrings-rs comparison uses original gates and truncation after each operation in both systems to match its gate-level coefficient decisions.
Other comparisons use the authors' SPD, MonoProp, PauliPropagation.jl, and ppvm implementations, Qiskit's Pauli propagation package, and our reproduction of ORQA.
The external implementations other than paulistrings-rs receive the exported Clifford-frame operation list.
MonoProp has a different truncation rule and requires separate error measurements for a comparison at equal accuracy.
The ORQA results describe our reproduction rather than the authors' implementation.
The default \SectorQ truncation interval is ten operations, with gate-level truncation used for the matched paulistrings-rs comparison.

\subsection{Performance Compared with Existing Systems}
The comparisons cover one node, seven circuits on 16 GPUs, and three circuits on 128 GPUs.
The paulistrings-rs experiment uses a separate matched truncation schedule.

\noindent\textbf{Single-Node Performance.}
Figure~\ref{fig:single} compares propagation time and completed cutoff ranges on one node.
Across circuits that fit one GPU, the ratio of the \cuPP time to the \SectorQ time is 0.6-1.5.
SPD is up to 2.9 times faster on one GPU, but its kicked Ising run exhausts memory at cutoff $10^{-5}$.
At their shared cutoff $3\times10^{-5}$, \SectorQ on four GPUs takes 5.25~s and distributed \cuPP takes 44.3~s.
\SectorQ also completes $3\times10^{-6}$ in 436~s.
At the largest shared settings across circuits, the four-GPU implementation is 108-269 times faster than PauliPropagation.jl, 157-755 than ppvm, and over 1,800 than single-threaded Qiskit.
The ORQA reproduction is 18-71 times slower on the same node.
MonoProp remains faster on several small settings below two seconds, including kicked Ising at $10^{-3}$.
The small-expansion measurements provide no consistent advantage for \SectorQ.
Its four-GPU execution completes larger expansions after several single-device implementations exhaust memory.

\begin{figure}[t]
\centering
\panel{ev_single_legend}\\[-2pt]
\panel{ev_single_a}
\caption{Simulation time on one node (4 GPUs)}
\label{fig:single}
\end{figure}

\noindent\textbf{Seven Circuits on 16 GPUs.}
Figure~\ref{fig:circuits} compares the seven circuits on four nodes.
At the largest setting completed by the original-circuit \cuPP baseline for each circuit, \SectorQ is 4.4-9.7 times faster.
Deferred deduplication remains 2.2-4.6 times slower at the shared settings and exhausts memory in two circuits at the primary baseline's largest completed setting.
\SectorQ completes lower cutoffs in all seven circuits, including TFIM and XXZ inputs without removable Clifford gates.
The TFIM and XXZ results demonstrate a reduction in propagation time even without removable Clifford operations.
For context, total propagation time divided by the maximum boundary count is 17-130~ns per term for \SectorQ above 100 million terms.
The primary baseline records 908-39,000~ns per term at cutoff $10^{-4}$.
These values describe amortized propagation time at different operator sizes, not a matched-size speedup or the cost of one gate per term.
% The ev_cost panels normalize these same time measurements by boundary count.
% Their summary is integrated here rather than repeating a seven-panel figure.

\begin{figure}[t]
\centering
\panel{ev_circuits_legend}\\[-2pt]
\subfloat[Kicked Ising]{\panel{ev_circuits_a}}\hfill
\subfloat[TFIM heavy-hex]{\panel{ev_circuits_b}}\hfill
\subfloat[Noisy kicked Ising]{\panel{ev_circuits_c}}\hfill
\subfloat[TFIM $11{\times}11$]{\panel{ev_circuits_d}}\\
\subfloat[XXZ $8{\times}8$]{\panel{ev_circuits_e}}\hfill
\subfloat[QAOA]{\panel{ev_circuits_f}}\hfill
\subfloat[Ansatz]{\panel{ev_circuits_g}}
\caption{Simulation time of seven circuits on 16 GPUs}
\label{fig:circuits}
\end{figure}

\noindent\textbf{Three Circuits on 128 GPUs.}
Figure~\ref{fig:circuits128} compares three circuits on 32 nodes using the original-circuit \cuPP baseline.
Kicked Ising and its noisy variant use those original inputs, while the TFIM input is identical with and without a Clifford frame.
At the largest setting completed by the primary baseline in each sweep, \SectorQ is 7.4-13.4 times faster.
For kicked Ising at cutoff $5\times10^{-6}$, the times are about 7.6 and 101.6~s.
The CPU comparisons use the same 32 nodes.
MonoProp remains competitive at the smallest settings, subject to its different truncation rule.
At the largest shared cutoffs completed by the ORQA reproduction, \SectorQ is 19.2-21.7 times faster.
These cutoffs are $3\times10^{-6}$ for kicked Ising and $2\times10^{-6}$ for TFIM and noisy kicked Ising.
The next ORQA setting exhausts memory for kicked Ising and its noisy variant, while the next TFIM setting is excluded by the time estimate.

\begin{figure}[t]
\centering
\panel{ev_circuits128_legend}\\[-2pt]
\subfloat[Kicked Ising]{\panel{ev_circuits128_a}}\hfill
\subfloat[TFIM heavy-hex]{\panel{ev_circuits128_b}}\hfill
\subfloat[Noisy kicked Ising]{\panel{ev_circuits128_c}}
\caption{Simulation time of three circuits on 128 GPUs}
\label{fig:circuits128}
\end{figure}

\noindent\textbf{Fixed-Partition Implementation with Matched Truncation.}
Figure~\ref{fig:prs} compares \SectorQ with paulistrings-rs on 16 GPUs using truncation after each operation and the same original circuit.
The fixed partition uses the library's seeded rows.
\SectorQ is 1.6-5.4 times faster over all 30 shared settings and 1.6-4.1 times faster at the largest shared setting of each circuit.
At the next tested cutoff in five circuits, paulistrings-rs exhausts GPU memory while \SectorQ completes in 29-880~s.
Times at small configurations include any lazy GPU kernel compilation within the propagation call.
This comparison includes differences in arithmetic kernels, storage, and partition management under a common truncation schedule.
The same-engine ablation in Section~\ref{sec:design-eval} isolates row recomputation, while Section~\ref{sec:accuracy} verifies the numerical agreement of the matched schedules.
% TBD: paulistrings-rs 32/64 GPU results from jobs 59769001 and 59769002.

\begin{figure}[t]
\centering
\panel{ev_prs_legend}\\[-2pt]
\subfloat[Kicked Ising]{\panel{ev_prs_a}}\hfill
\subfloat[TFIM heavy-hex]{\panel{ev_prs_b}}\hfill
\subfloat[Noisy kicked Ising]{\panel{ev_prs_c}}\hfill
\subfloat[TFIM $11{\times}11$]{\panel{ev_prs_d}}\\
\subfloat[XXZ $8{\times}8$]{\panel{ev_prs_e}}\hfill
\subfloat[QAOA]{\panel{ev_prs_f}}\hfill
\subfloat[Ansatz]{\panel{ev_prs_g}}
\caption{\SectorQ and fixed linear partitions on 16 GPUs}
\label{fig:prs}
\end{figure}

\subsection{Scalability and Capacity}
We vary the GPU count, circuit dimensions, and retained operator size to distinguish execution scaling from capacity.

\noindent\textbf{Strong Scalability.}
For fixed large problems, \SectorQ achieves speedups of 81-166 on 256 GPUs and 83-180 on 512 GPUs relative to one GPU.
The representative settings are kicked Ising at cutoff $3\times10^{-6}$ and heavy-hex TFIM and noisy kicked Ising at $5\times10^{-6}$.
Figure~\ref{fig:scaleout} also shows the kicked Ising cutoff ladder from one to 512 GPUs.
From 32 GPUs onward, its capacity measurements at cutoffs $10^{-6}$ and below use higher sector caps and, on 128 GPUs onward, smaller chunks.
The default cap of 256 sectors limits additional speedup beyond 256 GPUs.
At kicked Ising cutoff $3\times10^{-6}$, time decreases from 10.48 to 9.67~s.
The trace contains at most 256 nonempty sectors and no GPU splits, leaving fewer concurrent tasks than the 512 available workers.
On 512 GPUs, TFIM at cutoff $2\times10^{-6}$ takes 20.7~s with \SectorQ and 201.2~s with original-circuit \cuPP.
The next \cuPP setting, $10^{-6}$, exhausts GPU memory.

\begin{figure}[t]
\centering
\panel{ev_scaleout_legend}\\[-2pt]
\panel{ev_scaleout_a}
\caption{Kicked Ising propagation from 1 to 512 GPUs}
\label{fig:scaleout}
\end{figure}

\noindent\textbf{Weak Scalability.}
Figure~\ref{fig:weak} increases the boundary term count with the GPU count by lowering the cutoff.
Efficiency is boundary-term throughput per GPU relative to the corresponding one-GPU run, accounting for the counts actually produced.
For the large load, \SectorQ retains 80-88\% of this throughput on 16 GPUs and 72-80\% on 64 GPUs across the three circuits.
The range falls to 23-38\% on 512 GPUs under the default sector cap.
The measured large-load \cuPP configurations exhaust GPU memory, whereas the small load permits time comparisons between both systems.
On 16 GPUs, \SectorQ is 8.1-12 times faster than the original-circuit baseline for the small load.
On the same 64 nodes as the 256-GPU configuration, the ORQA reproduction takes 15.8-19.1 times longer for the small load.
MonoProp takes 3.6 and 5.0 times longer for kicked Ising and TFIM, respectively, under its different truncation rule.
The 512-GPU CPU and weak-load baseline configurations remain unexecuted.
The retained population increases with GPU count, but throughput per GPU decreases at the largest configurations under the default sector cap.
% Remaining: 512-GPU CPU and weak-load baseline points need a new allocation.

\begin{figure}[t]
\centering
\panel{ev_weak_legend}\\[-2pt]
\subfloat[Kicked Ising, small load]{\panel{ev_weak_a}}\hfill
\subfloat[TFIM heavy-hex, small load]{\panel{ev_weak_b}}\hfill
\subfloat[Noisy kicked Ising, small load]{\panel{ev_weak_c}}\\
\subfloat[Kicked Ising, large load]{\panel{ev_weak_d}}\hfill
\subfloat[TFIM heavy-hex, large load]{\panel{ev_weak_e}}\hfill
\subfloat[Noisy kicked Ising, large load]{\panel{ev_weak_f}}
\caption{Weak scaling from 1 to 512 GPUs}
\label{fig:weak}
\end{figure}

\noindent\textbf{Circuit Size and Depth.}
Figure~\ref{fig:size} varies the TFIM lattice size and kicked Ising depth on 16 GPUs.
The lattice sweep covers 36-1,024 qubits with total magnetization as the observable, changing string length, graph size, and observable size together.
At 1,024 qubits on a $32\times32$ lattice, \SectorQ completes in 90.4~s and is 4.1 times faster than original-circuit \cuPP.
The sweep ends at the library's supported limit of 1,024 qubits.
The depth sweep keeps the kicked Ising graph fixed and increases its Trotter steps.
At 60 steps, \SectorQ takes 160~s and is 50 times faster than the primary baseline.
The lower execution time persists across changes in string length, interaction graph, and propagation depth.
The term counts on the horizontal axes also vary across these configurations because retained operator size depends on the resulting branching and truncation.

\begin{figure}[t]
\centering
\panel{ev_size_legend}\\[-2pt]
\subfloat[Qubits, cutoff $10^{-3}$]{\panel{ev_size_a}}\hfill
\subfloat[Qubits, cutoff $10^{-4}$]{\panel{ev_size_b}}\hfill
\subfloat[Trotter steps]{\panel{ev_size_c}}
\caption{Simulation time by circuit size and depth}
\label{fig:size}
\end{figure}

\noindent\textbf{Largest Operator and Host Memory.}
The kicked Ising capacity sweep reaches 448.2 billion boundary terms on 512 GPUs at cutoff $1.5\times10^{-7}$, completing in 3,413~s.
The 17.9~TB expansion is 140 times larger than the largest completed original-circuit \cuPP setting on the same GPUs, 3.2 billion terms at $2\times10^{-6}$.
The host chunk pool occupies 21,079~GiB across 128 nodes, 64\% of installed host memory.
This allocation includes reusable chunks and transient storage and excludes memory outside the chunk pool.
On 128 GPUs, \SectorQ completes 117.5 billion terms.
On 256 GPUs, the completed endpoint is 257.3 billion terms at cutoff $2\times10^{-7}$ in 3,382~s.
Its host chunk pool occupies 11,490~GiB across 64 nodes, or 70\% of installed host memory.
Both CPU baselines on those 64 nodes complete kicked Ising at $2\times10^{-6}$ and exhaust memory at $10^{-6}$.
Configurations rejected from estimated storage requirements remain distinct from measured failures.
% Remaining: the 256-GPU ORQA TFIM/noise ladders and 512-GPU CPU reach are incomplete.

\subsection{Contribution of the Design Choices}
\label{sec:design-eval}
Controlled engine variants evaluate preprocessing and window-dependent rows.
Time counters quantify computation, copies, exchange, and worker idle time across GPU counts.

\noindent\textbf{Preprocessing and Window Sectors.}
Figure~\ref{fig:staircase} measures successive changes on kicked Ising at cutoff $10^{-5}$ and 16 GPUs.
The original-circuit \cuPP baseline takes 95.0~s.
A Clifford-frame input reduces time to 54.3~s, and deferred global deduplication reduces it to 48.8~s.
Window-sector execution reaches 12.2~s with the same underlying library operations.
The reduction from 48.8 to 12.2~s remains after Clifford preprocessing and deferred deduplication, identifying an additional benefit from sector execution.
The sequence changes the surrounding execution structure while retaining the library's single-GPU arithmetic.
The following fixed-row experiment compares partition rules within that sector engine.
Its recomputed variant takes 11.9~s at the same cutoff in a separate allocation.
The 11.9 and 12.2~s measurements are separate executions of the same design, and their difference is not attributed to another mechanism.

\begin{figure}[t]
\centering
\panel{ev_staircase_legend}\\[-2pt]
\panel{ev_staircase_a}
\caption{Effect of preprocessing and sector execution on 16 GPUs}
\label{fig:staircase}
\end{figure}

\noindent\textbf{Fixed versus Recomputed Rows.}
Figure~\ref{fig:fixedrows} isolates row recomputation with the same Clifford frame, GPU operations, and truncation schedule.
The fixed variant selects rows at the first partitioned window and retains them, ending a window when a rotation changes those rows.
The recomputed variant selects rows for each window.
Fixed rows are 2.4-17.0 times slower over 25 successful pairs on 16 GPUs and 1.3-14.8 times slower over 22 pairs on four GPUs.
All pairs have identical printed expectation values.
Several fixed-row runs exhaust host memory at lower cutoffs that the recomputed variant completes.
Noisy kicked Ising is excluded because the fixed variant does not support its explicit Clifford operations between rotations.
Table~\ref{tab:fixedrows} connects the time difference to window count and traffic.
For kicked Ising at $10^{-5}$, recomputation reduces the window count from 878 to 20 and inter-node traffic from 698.5 to 45.2~GB.
Fewer boundaries reduce the frequency of sector redistribution and the associated host-device transfers.

\begin{figure}[t]
\centering
\panel{ev_fixed_rows_legend}\\[-2pt]
\subfloat[Kicked Ising]{\panel{ev_fixed_rows_a}}\hfill
\subfloat[TFIM heavy-hex]{\panel{ev_fixed_rows_b}}\hfill
\subfloat[TFIM $11\times11$]{\panel{ev_fixed_rows_c}}\\
\subfloat[XXZ $8\times8$]{\panel{ev_fixed_rows_d}}\hfill
\subfloat[QAOA]{\panel{ev_fixed_rows_e}}\hfill
\subfloat[Ansatz]{\panel{ev_fixed_rows_f}}
\caption{Fixed and recomputed rows on 16 GPUs, with fixed-to-recomputed time ratios}
\label{fig:fixedrows}
\end{figure}

\begin{table}[t]
\caption{Windows and inter-node traffic at cutoff $10^{-5}$ on 16 GPUs}
\label{tab:fixedrows}
\centering\small
\begin{tabular}{@{}lrrrr@{}}
\toprule
& \multicolumn{2}{c}{\textbf{Windows}} & \multicolumn{2}{c}{\textbf{Traffic (GB)}} \\
\textbf{Circuit} & \textbf{Recomputed} & \textbf{Fixed} & \textbf{Recomputed} & \textbf{Fixed} \\
\midrule
Kicked Ising & 20 & 878 & 45.2 & 698.5 \\
TFIM heavy-hex & 45 & 2,241 & 79.0 & 342.2 \\
TFIM $11\times11$ & 24 & 1,518 & 60.0 & 625.9 \\
XXZ $8\times8$ & 23 & 1,155 & 44.8 & 2,369.9 \\
QAOA & 8 & 198 & 8.2 & 34.0 \\
Ansatz & 34 & 1,392 & 8.3 & 107.7 \\
\bottomrule
\end{tabular}
\end{table}

\noindent\textbf{Balanced Selection and Overlap.}
Figure~\ref{fig:resector-ablation} compares candidate selection and overlap within one allocation on 16 GPUs.
Panels (a) and (b) contain 30 circuit-cutoff pairs, with one measurement per variant and a fixed random seed.
Imbalance is the largest sector divided by the mean nonempty sector size, maximized over windows.
Eight candidates reduce this maximum in 28 pairs, with its median decreasing from 1.069 to 1.029.
The median one-to-eight-candidate time ratio is 1.00.
Median row-selection time increases from 0.113 to 0.459~s, limiting the overall time benefit.
Panel (c) compares overlap enabled and disabled for three circuits and two cutoffs.
Markers show ratios of mean times over three repetitions, and whiskers span the three paired ratios.
Disabling overlap increases mean wall time by 3-47\% across all six settings, with a geometric-mean increase of 21\%.
All 48 pairs print identical expectation values.
Candidate competition improves balance without a consistent independent speedup, and overlap reduces the exposed selection cost.

\noindent\textbf{Time Breakdown across GPU Counts.}
Figure~\ref{fig:runtime} uses cutoff $10^{-4}$ as a small reference and increases the other workloads with GPU count.
Within each circuit, the middle points keep boundary terms per GPU within a factor of 1.5.
The 64- and 256-GPU capacity points use sector caps of 1,024 and 4,096, with 65,536-term chunks on 256 GPUs.
The capacity configurations are marked with asterisks, and parentheses denote boundary term counts.
Worker times are averaged over GPUs, exchange and synchronization times over nodes, and each is normalized by wall time.
The Other component is the residual wall time after these recorded components.
The library-call component measures elapsed call time and does not represent sampled device utilization.
All 27 profiles are complete.
The largest 64-GPU points contain 14.9-78.5 billion boundary terms, and the largest 256-GPU points contain 37.3-212.6 billion.
Average worker idle time is below 1\% at all six capacity points.
The larger inputs sustain worker activity, although exchange and transfer costs remain.

\noindent\textbf{Concurrency and Exchange Costs.}
The small reference spends 37-52\% of wall time in average worker idle time across the nine profiles.
At the largest 16-GPU settings, library operations occupy 76-89\% of wall time and inter-node exchange occupies 3-7\%.
At the largest 256-GPU points, exchange occupies 10-35\%, with TFIM at the upper end.
The original-circuit \cuPP baseline separately records 39-44\% of its large-setting time inside MPI calls.
The separate 512-GPU measurements show that average worker idle time rises from 9\% to 40\% of wall time relative to 256 GPUs at kicked Ising cutoff $3\times10^{-6}$.
The largest capacity run uses a higher sector cap and has a different composition, with 46\% in library calls, 26\% in exchange, and 9\% in GPU-host copies.
Its exchange time is concentrated in data transfer and request posting, with under two seconds spent exchanging counts.
The fixed-input trace identifies limited task concurrency under the default cap, whereas the larger capacity run has a substantial exchange component despite low worker idle time.

\begin{figure}[tp]
\centering
\panel{ev_runtime_legend}\\[-2pt]
\subfloat[Kicked Ising, 16 GPUs]{\panel{ev_runtime_a}}\hfill
\subfloat[TFIM heavy-hex, 16 GPUs]{\panel{ev_runtime_b}}\hfill
\subfloat[Noisy kicked Ising, 16 GPUs]{\panel{ev_runtime_c}}\\[4pt]
\subfloat[Kicked Ising, 64 GPUs]{\panel{ev_runtime_d}}\hfill
\subfloat[TFIM heavy-hex, 64 GPUs]{\panel{ev_runtime_e}}\hfill
\subfloat[Noisy kicked Ising, 64 GPUs]{\panel{ev_runtime_f}}\\[4pt]
\subfloat[Kicked Ising, 256 GPUs]{\panel{ev_runtime_g}}\hfill
\subfloat[TFIM heavy-hex, 256 GPUs]{\panel{ev_runtime_h}}\hfill
\subfloat[Noisy kicked Ising, 256 GPUs]{\panel{ev_runtime_i}}
\caption{Time shares for a fixed small reference and larger workloads by GPU count}
\label{fig:runtime}
\end{figure}

\begin{figure}[t]
\centering
\panel{ev_resector_legend}\\[-2pt]
\subfloat[Sector balance]{\panel{ev_resector_a}}\hfill
\subfloat[Candidate selection time effect]{\panel{ev_resector_b}}\hfill
\subfloat[Overlap time effect]{\panel{ev_resector_c}}
\caption{Candidate selection and overlap on 16 GPUs, with ratios above one favoring eight candidates or enabled overlap}
\label{fig:resector-ablation}
\end{figure}

\subsection{Numerical Agreement and Accuracy}
\label{sec:accuracy}
We compare implementations under matched truncation schedules and assess approximation error against an exact reference.

\noindent\textbf{Agreement across Implementations.}
Across 68 circuit-cutoff settings with at least two successful runs under a common truncation rule, the largest absolute expectation difference is $4.9\times10^{-11}$.
This comparison includes the single-GPU reference, distributed \cuPP, and external implementations with the same coefficient decisions.
The largest difference occurs for QAOA's edge-sum observable.
The separate gate-level truncation comparison with paulistrings-rs agrees within $4.7\times10^{-11}$ over 75 completed settings on one, four, and 16 GPUs.
\SectorQ also preserves the printed expectation value when sector count, GPU count, window rule, or device-memory budget changes.
This agreement supports the merge-locality property across the evaluated partition and resource configurations.
It measures consistency under the specified truncation schedule, while approximation error requires a reference without that truncation.

\noindent\textbf{Error against an Exact Reference.}
Figure~\ref{fig:accuracy} uses a 30-qubit TFIM on a $5\times6$ lattice with an exact state-vector reference, checked against untruncated Pauli propagation on smaller grids.
The four-step circuit reaches absolute error $7.2\times10^{-5}$ at cutoff $10^{-7}$, with \SectorQ on 16 GPUs completing in 29.2~s.
Implementations with the same truncation rule have the same error at a given cutoff, whereas MonoProp requires a separate error comparison.
The six-step circuit reaches time or memory limits at higher cutoffs.
Error is nonmonotonic at intermediate cutoffs in both sweeps.
The large kicked Ising workload has no exact reference at the evaluated scale, and its expectation values are also nonmonotonic as the cutoff decreases.
For example, the value increases from approximately 0.3217 at $10^{-6}$ to 0.3265 at an intermediate cutoff, then decreases to 0.3089 at $1.5\times10^{-7}$.
The larger completed expansions establish computational capacity, but convergence of the kicked Ising expectation remains unverified.

\begin{figure}[t]
\centering
\panel{ev_accuracy_legend}\\[-2pt]
\subfloat[4 Trotter steps]{\panel{ev_accuracy_a}}\hfill
\subfloat[6 Trotter steps]{\panel{ev_accuracy_b}}
\caption{Simulation time and error of a 30-qubit TFIM}
\label{fig:accuracy}
\end{figure}

\subsection{Parameter Sensitivity and Repeatability}
Figure~\ref{fig:robustness} reports the effect of sector and memory parameters and the variability of repeated execution times.

\noindent\textbf{Sector Granularity and Device Memory.}
At cutoff $10^{-5}$, a target sector size one quarter of the default increases time by 9-26\%, whereas twice the target reduces it by 2-9\% across the representative circuits.
One to four tasks per GPU differ by at most 9\%, and eight can increase time by up to 35\% on 16 GPUs.
The tradeoff is between additional library invocations with small sectors and working-buffer pressure with large sectors.
With a device-memory budget of 4~GB, execution time increases by 12-64\% relative to 36~GB, while finer sectors preserve the tested expectation values.
A larger task count can also increase total time.
A 512-GPU kicked Ising run at $3\times10^{-6}$ with a 4,096-sector cap and smaller host chunks takes 11.6~s, against the default's 9.7~s.
The larger cap and smaller chunks change together in this experiment, preventing attribution of the time increase to either parameter alone.

\noindent\textbf{Repeated Measurements.}
The repetition experiment contains 160 completed kicked Ising runs, with ten measurements per setting.
We define spread as the difference between the maximum and minimum times divided by the mean of all repetitions.
On 16 GPUs, spread is 1.4-7.2\% for \SectorQ and 3.4-7.6\% for original-circuit \cuPP.
All repetitions preserve the printed expectation value.
Each system's 16-GPU repetitions use one allocation, with different allocations for the two systems.
The reported spread measures variation within those allocations.
At cutoff $10^{-5}$, the ratio of ten-run means is 7.6, close to the main sweep's single-run ratio of 7.8.
Panels (d) and (e) show the repeated times across problem sizes and GPU counts, respectively.
The difference between the systems' mean times persists across the ten repetitions at this configuration.

\begin{figure}[tp]
\centering
\panel{ev_knobs_legend}\\[-2pt]
\subfloat[Target terms per sector]{\panel{ev_knobs_a}}\hfill
\subfloat[Sectors per GPU]{\panel{ev_knobs_b}}\hfill
\subfloat[GPU memory allowed]{\panel{ev_knobs_c}}\\[4pt]
\panel{ev_variability_legend}\\[-2pt]
\subfloat[Repetitions on 16 GPUs by problem size]{\panel{ev_variability_a}}\hfill
\subfloat[Repetitions at cutoff $3{\times}10^{-5}$ by GPU count]{\panel{ev_variability_b}}
\caption{Parameter sensitivity at cutoff $10^{-5}$ and ten repeated measurements of kicked Ising}
\label{fig:robustness}
\end{figure}
